
\documentclass{article}
\usepackage{colortbl}
\usepackage{makecell}
\usepackage{multirow}
\usepackage{supertabular}

\begin{document}

\newcounter{utterance}

\centering \large Interaction Transcript for game `eqbench', experiment `eq\_bench', episode 7 with reflexion{-}llama{-}3.1{-}8b{-}instruct.
\vspace{24pt}

{ \footnotesize  \setcounter{utterance}{1}
\setlength{\tabcolsep}{0pt}
\begin{supertabular}{c@{$\;$}|p{.15\linewidth}@{}p{.15\linewidth}p{.15\linewidth}p{.15\linewidth}p{.15\linewidth}p{.15\linewidth}}
    \# & \multicolumn{2}{c}{Player} && \multicolumn{2}{c}{Game Master} \\
    \hline

    \theutterance \stepcounter{utterance}  
    & & & \multicolumn{4}{p{0.6\linewidth}}{
        \cellcolor[rgb]{0.9,0.9,0.9}{
            \makecell[{{p{\linewidth}}}]{
                \texttt{\tiny{[P1$\langle$GM]}}
                \texttt{You have attempted this same task before and did not succeed. These are the plans you wrote after those attempts.} \\
\texttt{Trial \#1: New plan: I overestimated the intensity of the emotions and didn't consider the context of the dialogue. In this conversation, Mike is expressing concern for Elise's well{-}being, and Elise is becoming defensive. The correct interpretation is that Mike is trying to help, but Elise is pushing him away. Next time, I will consider the context and the tone of the conversation to assign more accurate scores. I will also pay attention to the specific words and phrases used by the characters to gauge their emotional responses. For example, Elise's phrase "Oh, so now you're my therapist?" suggests that she is feeling defensive and possibly resentful, which I should have picked up on.} \\
\\ 
\texttt{Trial \#2: New plan: I overestimated the intensity of the emotions and didn't consider the context of the dialogue. In this conversation, Mike is expressing concern for Elise's well{-}being, and Elise is becoming defensive. The correct interpretation is that Mike is trying to help, but Elise is pushing him away. Next time, I will consider the context and the tone of the conversation to assign more accurate scores. I will also pay attention to the specific words and phrases used by the characters to gauge their emotional responses. For example, Elise's phrase "Oh, so now you're my therapist?" suggests that she is feeling defensive and possibly resentful, which I should have picked up on.} \\
\\ 
\texttt{To improve my scores, I will pay closer attention to the following:} \\
\\ 
\texttt{{-} Elise's defensive tone and language, such as "Oh, so now you're my therapist?" and "Well, 'friend', I didn't ask for your worry."} \\
\texttt{{-} Mike's concern and attempt to help, which is evident in his words "You've been working too much, Elise" and "I'm your friend. And as your friend, I'm worried about you."} \\
\texttt{{-} The context of the conversation, which suggests that Elise is pushing Mike away and not receptive to his concern.} \\
\\ 
\texttt{With these factors in mind, I will reassess the emotions and assign new scores:} \\
\\ 
\texttt{Relief: 0} \\
\texttt{Resentment: 8} \\
\texttt{Surprise: 0} \\
\texttt{Frustration: 6} \\
\\ 
\texttt{Trial \#3: New plan: I overestimated the intensity of the emotions and didn't consider the context of the dialogue. In this conversation, Mike is expressing concern for Elise's well{-}being, and Elise is becoming defensive. The correct interpretation is that Mike is trying to help, but Elise is pushing him away. Next time, I will consider the context and the tone of the conversation to assign more accurate scores. I will also pay attention to the specific words and phrases used by the characters to gauge their emotional responses. For example, Elise's phrase "Oh, so now you're my therapist?" suggests that she is feeling defensive and possibly resentful, which I should have picked up on.} \\
\\ 
\texttt{However, upon re{-}examining my previous attempts, I realize that I made the same mistake twice. This suggests that I need to take a more nuanced approach to understanding the emotions expressed in the dialogue. To improve my scores, I will pay closer attention to the following:} \\
\\ 
\texttt{{-} Elise's defensive tone and language, such as "Oh, so now you're my therapist?" and "Well, 'friend', I didn't ask for your worry."} \\
\texttt{{-} Mike's concern and attempt to help, which is evident in his words "You've been working too much, Elise" and "I'm your friend. And as your friend, I'm worried about you."} \\
\texttt{{-} The context of the conversation, which suggests that Elise is pushing Mike away and not receptive to his concern.} \\
\\ 
\texttt{With these factors in mind, I will reassess the emotions and assign new scores:} \\
\\ 
\texttt{To better understand Elise's emotions, I will consider the following:} \\
\\ 
\texttt{{-} Her phrase "And you've been reading too much" could be seen as a dismissive or defensive response, which might indicate a higher level of resentment.} \\
\texttt{{-} Her phrase "Well, 'friend', I didn't ask for your worry" suggests that she is feeling defensive and possibly resentful, which I should have picked up on.} \\
\\ 
\texttt{Considering these factors, I will reassess the emotions and assign new scores:} \\
\\ 
\texttt{Relief: 0} \\
\texttt{Resentment: 9} \\
\texttt{Surprise: 0} \\
\texttt{Frustration: 4} \\
\\ 
\\ 
\\ 
\texttt{Your task is to predict the likely emotional responses of a character in this dialogue:} \\
\\ 
\texttt{Mike: You've been working too much, Elise.} \\
\texttt{Elise: And you've been reading too much. We all have our vices, Mike.} \\
\texttt{Mike: It's not a vice, it's a concern. You're burning yourself out.} \\
\texttt{Elise: Oh, so now you're my therapist? } \\
\texttt{Mike: No, I'm your friend. And as your friend, I'm worried about you.} \\
\texttt{Elise: Well, 'friend', I didn't ask for your worry.} \\
\texttt{{[}End dialogue{]}} \\
\\ 
\texttt{At the end of this dialogue, Mike would feel...} \\
\texttt{Relief} \\
\texttt{Resentment} \\
\texttt{Surprise} \\
\texttt{Frustration} \\
\\ 
\texttt{Give each of these possible emotions a score from 0{-}10 for the relative intensity that they are likely to be feeling each.} \\
\\ 
\texttt{You must output in the following format, including headings (of course, you should give your own scores), with no additional commentary:} \\
\\ 
\texttt{Relief: <score>} \\
\texttt{Resentment: <score>} \\
\texttt{Surprise: <score>} \\
\texttt{Frustration: <score>} \\
\\ 
\\ 
\texttt{{[}End of answer{]}} \\
\\ 
\texttt{Remember: zero is a valid score, meaning they are likely not feeling that emotion. You must score at least one emotion > 0.} \\
\\ 
\texttt{Your answer:} \\
            }
        }
    }
     \\ \\

    \theutterance \stepcounter{utterance}  
    & \multicolumn{4}{p{0.6\linewidth}}{
        \cellcolor[rgb]{0.561,0.737,0.561}{
            \makecell[{{p{\linewidth}}}]{
                \texttt{\tiny{[P1$\rangle$GM]}}
                \texttt{Relief: 0} \\
\texttt{Resentment: 8} \\
\texttt{Surprise: 0} \\
\texttt{Frustration: 6} \\
            }
        }
    }
    & & \\ \\

    \theutterance \stepcounter{utterance}  
    & & & \multicolumn{2}{p{0.3\linewidth}}{
        \cellcolor[rgb]{0.9,0.9,0.9}{
            \makecell[{{p{\linewidth}}}]{
                \texttt{\tiny{[GM$|$GM]}}
                \texttt{Relief: 0} \\
\texttt{Resentment: 8} \\
\texttt{Surprise: 0} \\
\texttt{Frustration: 6} \\
            }
        }
    }
    & & \\ \\

    \theutterance \stepcounter{utterance}  
    & & & \multicolumn{2}{p{0.3\linewidth}}{
        \cellcolor[rgb]{0.9,0.9,0.9}{
            \makecell[{{p{\linewidth}}}]{
                \texttt{\tiny{[GM$|$GM]}}
                \texttt{Relief: 0} \\
\texttt{Resentment: 3} \\
\texttt{Surprise: 2} \\
\texttt{Frustration: 8} \\
            }
        }
    }
    & & \\ \\

    \theutterance \stepcounter{utterance}  
    & & & \multicolumn{2}{p{0.3\linewidth}}{
        \cellcolor[rgb]{0.9,0.9,0.9}{
            \makecell[{{p{\linewidth}}}]{
                \texttt{\tiny{[GM$|$GM]}}
                \texttt{game\_result = LOSE} \\
            }
        }
    }
    & & \\ \\

    \theutterance \stepcounter{utterance}  
    & & \multicolumn{4}{p{0.6\linewidth}}{
        \cellcolor[rgb]{0.95,0.95,0.95}{
            \makecell[{{p{\linewidth}}}]{
                \texttt{\tiny{[P1$\leftrightarrow$P]}}
                \texttt{Post{-}episode Reflexion (not sent to GM):} \\
\texttt{New plan: I overestimated the intensity of the emotions and didn't consider the context of the dialogue. In this conversation, Mike is expressing concern for Elise's well{-}being, and Elise is becoming defensive. The correct interpretation is that Mike is trying to help, but Elise is pushing him away. Next time, I will consider the context and the tone of the conversation to assign more accurate scores. I will also pay attention to the specific words and phrases used by the characters to gauge their emotional responses. For example, Elise's phrase "Oh, so now you're my therapist?" suggests that she is feeling defensive and possibly resentful, which I should have picked up on.} \\
\\ 
\texttt{However, upon re{-}examining my previous attempts, I realize that I made the same mistake twice. This suggests that I need to take a more nuanced approach to understanding the emotions expressed in the dialogue. To improve my scores, I will pay closer attention to the following:} \\
\\ 
\texttt{{-} Elise's defensive tone and language, such as "Oh, so now you're my therapist?" and "Well, 'friend', I didn't ask for your worry."} \\
\texttt{{-} Mike's concern and attempt to help, which is evident in his words "You've been working too much, Elise" and "I'm your friend. And as your friend, I'm worried about you."} \\
\texttt{{-} The context of the conversation, which suggests that Elise is pushing Mike away and not receptive to his concern.} \\
\\ 
\texttt{With these factors in mind, I will reassess the emotions and assign new scores:} \\
\\ 
\texttt{To better understand Elise's emotions, I will consider the following:} \\
\\ 
\texttt{{-} Her phrase "And you've been reading too much" could be seen as a dismissive or defensive response, which might indicate a higher level of resentment.} \\
\texttt{{-} Her phrase "Well, 'friend', I didn't ask for your worry" suggests that she is feeling defensive and possibly resentful, which I should have picked up on.} \\
\\ 
\texttt{Considering these factors, I will reassess the emotions and assign new scores:} \\
\\ 
\texttt{Relief: 0} \\
\texttt{Resentment: 9} \\
\texttt{Surprise: 0} \\
\texttt{Frustration: 4} \\
            }
        }
    }
    & \\ \\

\end{supertabular}
}

\end{document}
